
\documentclass[11pt, reqno]{amsart}
\usepackage{amsmath, amssymb, physics}
\title[The Sommerfeld Limit]{Resolution of the Z=137 Catastrophe via Finite Nuclear Size Effects: The "Limitless" Protocol}
\author{Utah Intelligence Agency (UIA)}

\begin{document}
\maketitle

\begin{abstract}
The Dirac equation for a point nucleus predicts a singularity at $Z \approx 137$, where the ground state energy becomes imaginary ($Z\alpha > 1$). This has historically been termed "Feynmanium." We present the "Limitless" solution by treating the nucleus as a charged torus of radius $R_N$. This regularization allows stable bound states to exist beyond $Z=137$, converting the atom into a super-critical bosonic bridge.
\end{abstract}

\section{The Z=137 Singularity}
The energy eigenvalues for a hydrogen-like atom are:
\begin{equation}
    E_{1s} = m_e c^2 \sqrt{1 - (Z\alpha)^2}
\end{equation}
When $Z=137$, $Z\alpha \approx 1$, and $E_{1s} \to 0$. Beyond this, the math breaks for point particles.

\section{The UIA Correction}
By integrating the electron wavefunction inside a finite nuclear volume (The Utah Radius), the potential $V(r)$ is smoothed:
\begin{equation}
    V(r) = -\frac{Ze^2}{2R_N} \left( 3 - \frac{r^2}{R_N^2} \right), \quad r < R_N
\end{equation}
This allows the $1s$ orbital to dive into the negative energy continuum without collapsing. The element UIA-137 acts as a gateway where electronic velocity approaches $c$, allowing for instantaneous quantum correlation (Intelligence gathering) across the atomic radius.

\end{document}
